# Dehn Twists on the Deformation Space of a Quadratic Rational Map with a Critical Point of Period Four

```latex
\documentclass[11pt]{amsart}
\usepackage{amssymb,amsthm,amsmath,amscd}
\usepackage{hyperref}

\theoremstyle{plain}
\newtheorem{theorem}{Theorem}[section]
\newtheorem{proposition}[theorem]{Proposition}
\newtheorem{lemma}[theorem]{Lemma}
\newtheorem{corollary}[theorem]{Corollary}
\theoremstyle{definition}
\newtheorem{definition}[theorem]{Definition}
\newtheorem{remark}[theorem]{Remark}
\newtheorem{example}[theorem]{Example}

\newcommand{\Sph}{\mathbb{S}^2}
\newcommand{\C}{\mathbb{C}}
\newcommand{\R}{\mathbb{R}}
\newcommand{\Z}{\mathbb{Z}}
\newcommand{\H}{\mathbb{H}}
\newcommand{\Teich}{\mathrm{Teich}}
\newcommand{\Mod}{\mathrm{Mod}}
\newcommand{\Homeo}{\mathrm{Homeo}}
\newcommand{\QC}{\mathrm{QC}}
\newcommand{\PSL}{\mathrm{PSL}}
\newcommand{\D}{\mathcal{D}}
\newcommand{\M}{\mathcal{M}}
\newcommand{\S}{\mathcal{S}}
\newcommand{\P}{\mathcal{P}}
\newcommand{\Crit}{\mathrm{Crit}}
\newcommand{\id}{\mathrm{id}}
\newcommand{\Spec}{\mathrm{Spec}}

\begin{document}

\title{Dehn twists and the deformation space of a quadratic rational map\\
with a critical point of period four}

\author{}
\date{}

\begin{abstract}
Let $f$ be a quadratic rational map with a critical point $c_1$ of period
$4$ whose cycle is $A=\{c_1,f(c_1),f^{\circ 2}(c_1),f^{\circ 3}(c_1)\}$, and
whose second critical value $v=f(c_2)$ does not lie on the cycle.  For
$B=A\cup\{v\}$ we consider Epstein's deformation space
$\D_f=\{\tau\in\Teich(\Sph,B)\,:\, f^{*}\tau=i_{*}\tau\}$.  We prove that the
Dehn twist $T_\gamma$ about a simple closed curve $\gamma\subset(\Sph,B)$
separating $\{c_1,v\}$ from $\{f(c_1),f^{\circ 2}(c_1),f^{\circ 3}(c_1)\}$
has the property that the square $T_\gamma^{\,2}$ preserves $\D_f$, and that
the induced action of $T_\gamma^{\,2}$ on the set $\pi_0(\D_f)$ of connected
components of $\D_f$ has infinite orbits.  In particular, $\D_f$ has
infinitely many connected components.
\end{abstract}

\maketitle

\section{Introduction and statement of the main theorem}\label{s:intro}

\subsection{Setup}

Let $f:\Sph\to\Sph$ be a map of degree two, i.e.\ a quadratic rational map
once a conformal identification $\Sph\cong\widehat\C$ has been chosen.  Let
$c_1,c_2$ be its critical points, let $V(f)=\{f(c_1),f(c_2)\}$ be its set of
critical values, and suppose:

\begin{itemize}
\item $c_1$ is periodic of \emph{exact} period $4$; write
$$A=\{c_1,\ p_1,\ p_2,\ p_3\}=\{c_1,\ f(c_1),\ f^{\circ 2}(c_1),\
f^{\circ 3}(c_1)\},$$
so that $f(A)=A$ and $f|_A$ is the cyclic permutation
$c_1\mapsto p_1\mapsto p_2\mapsto p_3\mapsto c_1$;
\item the second critical value $v:=f(c_2)$ satisfies $v\notin A$.
\end{itemize}
Set $B:=A\cup\{v\}$.  Note that $f(A)\cup V(f)\cup A\subseteq B$: indeed
$f(A)=A$ and $V(f)=\{p_1,v\}\subseteq B$.  Hence, in the notation of the
statement of the problem, both the pullback map
$f^{*}:\Teich(\Sph,B)\to\Teich(\Sph,A)$ and the forgetful map
$i_{*}:\Teich(\Sph,B)\to\Teich(\Sph,A)$ are defined, and so is the
\emph{deformation space}
$$\D_f:=\{\tau\in\Teich(\Sph,B)\ :\ f^{*}\tau=i_{*}\tau\}.$$

\subsection{Main theorem}

For a finite set $X\subset\Sph$ with $|X|\ge 3$, a simple closed curve in
$(\Sph,X)$ is a Jordan curve $\gamma\subset\Sph\setminus X$; it is
\emph{essential} if the two components of $\Sph\setminus\gamma$ each contain
at least two points of $X$.  Essential curves are classified, up to isotopy
rel.\ $X$, by the partition of $X$ induced by the two sides.

\begin{theorem}[Main Theorem]\label{thm:main}
Let $f$ be a quadratic rational map with critical points $c_1,c_2$, where
$c_1$ has exact period $4$ with cycle $A$, and $v=f(c_2)\notin A$; let
$B=A\cup\{v\}$ and let $\D_f\subset\Teich(\Sph,B)$ be the associated
deformation space.

Let $\gamma\subset(\Sph,B)$ be any essential simple closed curve such that
one of the two components of $\Sph\setminus\gamma$ is a disk $D$ with
$D\cap B=\{c_1,v\}$ (i.e.\ $\gamma$ ``encloses exactly the critical point
$c_1$ and the second critical value $v$'').  Then:
\begin{enumerate}
\item[(a)] the square of the Dehn twist $T_\gamma$ about $\gamma$ preserves
$\D_f$; that is, $T_\gamma^{\,2}\cdot\D_f=\D_f$ inside $\Teich(\Sph,B)$;
\item[(b)] the induced action of $T_\gamma^{\,2}$ on the set $\pi_0(\D_f)$
of connected components of $\D_f$ has infinite orbits.  In particular
$\D_f$ has infinitely many connected components.
\end{enumerate}
The same conclusions hold for the curves enclosing exactly $\{p_2,v\}$ or
exactly $\{p_3,v\}$.
\end{theorem}

\begin{remark}
(i) Part (a) is \emph{unconditional}: it is a purely topological
combinatorial statement about the branched cover
$f:(\Sph,f^{-1}B)\to(\Sph,B)$, proved in \S\ref{s:preserve} by an explicit
computation of the full inverse image of $\gamma$.

(ii) The reason a \emph{power} of the twist is needed is that
$f^{-1}(\gamma)$ is a \emph{single} curve mapping onto $\gamma$ with degree
$2$; a Dehn twist about $\gamma$ admits a lift through $f$ only after
squaring.  This is proved in Lemma~\ref{lem:lift}.  The factor $2$ matches
the fact, established in \S\ref{s:stratum}, that in the explicit
one--parameter family of maps with a marked $4$--cycle, the second critical
value is the \emph{square} of a rational function of the parameter, so that
one turn of the parameter around the corresponding cusp produces two turns
of $v$ around $c_1$.
\end{remark}

\subsection{Strategy and scheme of proof}

The proof of Theorem~\ref{thm:main} proceeds through the following steps.

\begin{enumerate}
\item[\S\ref{s:prelim}] Foundations: Teichm\"uller spaces of marked spheres,
mapping class groups, peripheral twists, the forgetful map and its
equivariance, and the point--pushing subgroup (Birman exact sequence).
\item[\S\ref{s:pullback}] The pullback map $f^{*}$ and the dynamical
characterisation of $\D_f$: a point $\tau=[\phi]$ of $\Teich(\Sph,B)$ lies
in $\D_f$ if and only if there exists a quasiconformal $u\sim_A\id$ such
that $\phi\circ f\circ u\circ\phi^{-1}$ is rational
(Proposition~\ref{prop:dynchar}).
\item[\S\ref{s:twists}] Lifting twists through the double cover: the exact
exponent formula $k/d_i$ for the lift of $T_\gamma^{k}$, the two
intertwining identities (Lemmas \ref{lem:intertwine1}, \ref{lem:intertwine2}),
and the \emph{Preservation Criterion} (Theorem~\ref{thm:crit}).
\item[\S\ref{s:stratum}] The explicit algebraic model of the stratum of
quadratic rational maps with a marked critical $4$--cycle: the family
$f_t$, the formula $v(t)=\bigl(\tfrac{t^2-t-1}{t^2+t-1}\bigr)^{2}$, the set
of cusps, and the braid monodromy table
(Lemma~\ref{lem:monodromytable}).
\item[\S\ref{s:preserve}] Proof of Theorem~\ref{thm:main}(a) by the
criterion of \S\ref{s:twists}, with the complete case analysis.
\item[\S\ref{s:components}] The component structure of $\D_f$: the
identification of the fibre of $\D_f\to\S$ over the base point with a
subgroup $Z\le\Mod(\Sph,B)$ (Proposition~\ref{prop:fibre}), the monodromy
group $M\subseteq Z$, the coset description
$\pi_0(\D_f)\cong Z/M$ (Proposition~\ref{prop:cosets}), the fact that
$Z$ contains a nonabelian free group generated by the squares of the three
pushes of $v$ (Proposition~\ref{prop:free}), and the deduction of
Theorem~\ref{thm:main}(b) from the class--monodromy input at the cusps
(Proposition~\ref{prop:input}), which we prove using the explicit model of
\S\ref{s:stratum} together with the standard theory of invariant curves for
critically preperiodic maps.
\end{enumerate}

Throughout, ``isotopy'' always means isotopy through orientation--preserving
homeomorphisms fixing each point of the relevant marked set, and
$\Homeo^+(\Sph,X)$ denotes the group of orientation--preserving
homeomorphisms of $\Sph$ fixing $X$ pointwise.

\section{Marked spheres, Teichm\"uller spaces, and mapping class groups}
\label{s:prelim}

\subsection{Teichm\"uller spaces}

Let $X\subset\Sph$ be a finite set.

\begin{definition}\label{def:teich}
$\Teich(\Sph,X)$ is the set of equivalence classes of quasiconformal
homeomorphisms $\phi:\Sph\to\widehat\C$, where
$\phi\sim_X\phi'$ if and only if there exist a M\"obius transformation
$M\in\PSL_2(\C)$ with $M(\phi(x))=\phi'(x)$ for all $x\in X$, and a
quasiconformal homeomorphism $h:\Sph\to\Sph$ with $h|_X=\id|_X$ and
$h$ isotopic to $\id$ rel.\ $X$, such that $\phi'=M\circ\phi\circ h$.
We write $[\phi]_X$ for the class of $\phi$.
\end{definition}

If $|X|\le 3$ then $\Teich(\Sph,X)$ is a point; if $|X|=n\ge4$ then
$\Teich(\Sph,X)$ is a complex manifold of dimension $n-3$, naturally
identified with the Teichm\"uller space of the $n$--punctured sphere.  We
only need:
\begin{equation}\label{eq:dim}
\dim_\C\Teich(\Sph,B)=2,\qquad \dim_\C\Teich(\Sph,A)=1 .
\end{equation}

\subsection{Mapping class groups and the pure braid group}

The \emph{pure mapping class group} of $(\Sph,X)$ is
$$\Mod(\Sph,X):=\Homeo^+(\Sph,X)/\bigl(\text{isotopy rel.\ }X\bigr).$$
It acts on $\Teich(\Sph,X)$ on the left by
\begin{equation}\label{eq:action}
g\cdot[\phi]_X:=[\phi\circ g^{-1}]_X,\qquad g\in\Mod(\Sph,X),
\end{equation}
which is well defined because pre--composition by an $X$--isotopy changes the
marking within its class, and post--composition by M\"obius maps is quotiented
out.  The group $\Mod(\Sph,X)$ is naturally the pure braid group of the
sphere on $|X|$ strands, and for $|X|\ge 4$ the action \eqref{eq:action} on
$\Teich(\Sph,X)$ is free and properly discontinuous.

\begin{lemma}[Normalization]\label{lem:normalization}
Suppose $|X|\ge3$ and $\phi,\phi'\in\QC$ satisfy $\phi'(x)=\phi(x)$ for all
$x\in X$.  Then $[\phi']_X=[\phi]_X$ if and only if $\phi'$ and $\phi$ are
isotopic rel.\ $X$ through quasiconformal homeomorphisms.  Consequently, if
$\pi:\Teich(\Sph,X)\to\M_X:=\Teich(\Sph,X)/\Mod(\Sph,X)$ is the quotient,
and $s\in\M_X$, then every $\tau\in\pi^{-1}(s)$ admits a unique
representative $\phi$ with $\phi|_X$ equal to any prescribed marking--compatible
embedding of $X$ into $\widehat\C$.
\end{lemma}

\begin{proof}
If $\phi'=M\phi h$ with $h\sim_X\id$ and $M(\phi(x))=\phi'(x)=\phi(x)$ for
all $x\in X$ and $|X|\ge3$, then $M$ fixes at least three points, so $M=\id$
and $\phi'=\phi h$ with $h\sim_X\id$, which is exactly an isotopy rel.\ $X$.
The last statement follows because $\M_X$ is the moduli of ordered
configurations: two markings with the same value of $\phi|_X$ lie over the
same $s$, and differ by a unique deck transformation, represented rel.\ $X$
by a unique isotopy class.
\end{proof}

\subsection{Curves, Dehn twists, and peripheral twists}

For an essential curve $\gamma$ in $(\Sph,X)$, the Dehn twist $T_\gamma\in
\Mod(\Sph,X)$ is the class of a homeomorphism supported in an annular
neighbourhood of $\gamma$, twisting once in the positive direction.  Two
essential curves in $(\Sph,X)$ are isotopic rel.\ $X$ if and only if they
induce the same partition of $X$.

\begin{lemma}[Peripheral twists are trivial]\label{lem:peripheral}
Let $X\subset\Sph$ be finite with $|X|\ge 1$ and let $\delta$ be a Jordan
curve bounding a disk $D$ with $D\cap X=\{x\}$ a single point.  Then
$T_\delta=\id$ in $\Mod(\Sph,X)$.
\end{lemma}

\begin{proof}
Realize $\delta$ as the circle $\{|z-x|=1\}$ in a coordinate $z$ with
$x=0$, and take the twist supported in the annulus
$\{\tfrac14\le|z|\le\tfrac12\}$, i.e.\ $T(z)=z\,e^{2\pi i\theta(|z|)}$ with
$\theta$ a smooth bump vanishing near $|\cdot|=\tfrac14,\tfrac12$ and
$\theta\le 1$.  For each $s\in[0,1]$ let $T_s$ be the twist supported in
$\{\tfrac{\epsilon(s)}2\le|z|\le\tfrac12\}$ with profile $\theta$, where
$\epsilon(0)=\tfrac12$ and $\epsilon(s)\to0$.  Then
$T_s(x)=x$ for all $x\in X$ (as $x=0$ lies inside the inner disk, where
$T_s=\id$) and $T_s\to\id$ uniformly as $s\to1$.  The family
$T_{s}$, extended continuously to $s=1$ by $T_1=\id$, is an isotopy rel.\
$X$ (each $T_s$ is a homeomorphism fixing $X$ pointwise, and the extension
is continuous).  Hence $T_\delta\sim_X\id$.
\end{proof}

\begin{corollary}\label{cor:peripheral-action}
If $\gamma$ bounds a disk whose intersection with $X$ has at most one
point, then $T_\gamma$ acts trivially on $\Teich(\Sph,X)$ and on curves
there is no essential twist; in particular, in the criterion of
Theorem~\ref{thm:crit} below, such curves contribute nothing.
\end{corollary}

\subsection{The forgetful map and its equivariance}

Let $X\subseteq Y$ be finite sets in $\Sph$.  The \emph{forgetful map}
$$i_*:\Teich(\Sph,Y)\longrightarrow\Teich(\Sph,X),\qquad
[\phi]_Y\longmapsto[\phi]_X$$
is well defined, because the equivalence relation defining
$[\,\cdot\,]_X$ is coarser than that defining $[\,\cdot\,]_Y$ (isotopy rel.\
$X$ is a weaker requirement than isotopy rel.\ $Y$, and the M\"obius
matching condition involves fewer points).  There is a group homomorphism
\begin{equation}\label{eq:forgethom}
\pi_1^{Y,X}:\Mod(\Sph,Y)\to\Mod(\Sph,X),\qquad
[h]_Y\mapsto [h]_X,
\end{equation}
which is well defined because an isotopy rel.\ $Y$ is also an isotopy rel.\
$X$.

\begin{lemma}[Equivariance of the forgetful map]\label{lem:intertwine1}
For every $g\in\Mod(\Sph,Y)$ and every $\tau\in\Teich(\Sph,Y)$,
$$i_*(g\cdot\tau)=\pi_1^{Y,X}(g)\cdot i_*(\tau).$$
In particular, if $\gamma$ is a curve in $(\Sph,Y)$ whose image
$\bar\gamma$ in $(\Sph,X)$ is non--essential (bounds at most one point of
$X$ on each side), then $i_*\circ T_\gamma=i_*$, by
Lemma~\ref{lem:peripheral}.
\end{lemma}

\begin{proof}
With the action \eqref{eq:action}, write $\tau=[\phi]_Y$ and let $h$
represent $g$.  Then
$i_*(g\cdot\tau)=[\phi\circ h^{-1}]_X
=[\phi]_X\circ$ (pre--composition by $[h^{-1}]_X$)
$=\pi_1^{Y,X}(g)\cdot[\phi]_X$.
\end{proof}

\subsection{Point pushing}

We shall use the classical Birman exact sequence in the following form.

\begin{lemma}[Birman; see \cite{Birman}, Ch.~4, and \cite{FM},
Thm.~4.6]\label{lem:birman}
Let $X\subset\Sph$, $|X|\ge3$, and $y\in\Sph\setminus X$.  Then the kernel
of the forgetful homomorphism
$\pi_1^{X\cup\{y\},\,X}:\Mod(\Sph,X\cup\{y\})\to\Mod(\Sph,X)$
is naturally isomorphic to $\pi_1(\Sph\setminus X)$: the element
represented by a loop $\beta$ based at $y$ maps to the class of any
homeomorphism obtained by pushing $y$ once along $\beta$.  Moreover, for a
small embedded loop $\beta_x$ winding once positively around
$x\in X$, the corresponding ``push'' is the Dehn twist about the curve
enclosing exactly $\{x,y\}$:
$$\mathrm{push}(\beta_x)=T_{\{x,y\}}.$$
\end{lemma}

\begin{corollary}\label{cor:pushfree}
Let $X=\{x_1,\dots,x_n\}$, $n\ge 4$.  The subgroup
$\Pi:=\ker\pi_1^{B,X}\cong\pi_1(\Sph\setminus X)$ is free of rank $n-1$,
freely generated by the pushes along $\beta_{x_1},\dots,\beta_{x_{n-1}}$.
The subgroup of $\Pi$ generated by the \emph{squares} of any $n-2$ of the
generators is free of rank $n-2$, with those squares as a free basis.
\end{corollary}

\begin{proof}
Freeness is the standard presentation of the fundamental group of a
punctured sphere.  For the second statement, let $F=F(a_1,\dots,a_{n-2})$ and
consider a nonempty reduced word $w$ in $a_1^{\pm2},\dots,a_{n-2}^{\pm2}$;
when expanded in the letters $a_i^{\pm1}$ it is a reduced nonempty word,
because at every junction of two consecutive factors the two boundary
letters carry the same sign of the same $a_i$, hence do not cancel.  So
there are no relations among the squares.
\end{proof}

\section{Branched covers and the pullback map}\label{s:pullback}

\subsection{Definition of the pullback}

Let $f:\Sph\to\Sph$ be a branched cover, let $A,B\subset\Sph$ be finite with
$f(A)\cup V(f)\cup A\subseteq B$ and $f(A)\subseteq A$ (in our situation
$f(A)=A$).  Fix once and for all a conformal identification
$\Sph\cong\widehat\C$, so that $f$ is a rational map and
$A,B\subseteq\widehat\C$ are honest configurations.

\begin{definition}\label{def:pullback}
Let $\tau=[\phi]_B\in\Teich(\Sph,B)$ and let
$\sigma_\phi:=\phi^{*}\sigma_0$ be the pulled--back complex structure
($\sigma_0$ the standard structure of $\widehat\C$).  Let
$f^{*}\sigma_\phi$ be the pullback of $\sigma_\phi$ under $f$, a complex
structure on $\Sph$ for which
$$\phi\circ f:\ (\Sph,\ f^{*}\sigma_\phi)\ \longrightarrow\ \widehat\C$$
is holomorphic.  Choose a quasiconformal homeomorphism
$\psi:\Sph\to\widehat\C$ with $\psi^{*}\sigma_0=f^{*}\sigma_\phi$, and set
$$f^{*}\tau:=[\psi]_A\in\Teich(\Sph,A),$$
where the marking is the one induced by the identity on the underlying set
$A\subset\Sph=\widehat\C$.
\end{definition}

\begin{lemma}[Well--definedness]\label{lem:welldefined}
Definition~\ref{def:pullback} is independent of the representative $\phi$
of $\tau$ and of the solution $\psi$.
\end{lemma}

\begin{proof}
Independence of $\psi$ is clear: two solutions differ by post--composition
with a M\"obius map, which does not change the class in $\Teich(\Sph,A)$
provided the marked points are matched; and any two normalized solutions
agree on the three points used for normalization, so the M\"obius map fixes
the marked triple and can be absorbed.

Suppose $\phi'=\phi h$ with $h\sim_B\id$ quasiconformal (the M\"obius factor
is harmless since $(M\phi)^{*}\sigma_0=\phi^{*}\sigma_0$).  Put
$u_t$, $t\in[0,1]$, the isotopy with $u_0=\id$, $u_1=h$, each $u_t$ fixing
$B$ pointwise.  Since $f(A)=A\subseteq B$, each composite $u_t\circ f$ fixes
$A$ pointwise, and $t\mapsto u_t\circ f$ is an isotopy of branched covers
through maps fixing $A$.  Consider the family of complex structures
$(u_t\circ f)^{*}\sigma_\phi$; solving the Beltrami equation continuously in
$t$ with a fixed normalization at three points of $A$ yields a continuous
family of markings $\psi_t$ with $\psi_t^*\sigma_0=(u_t\circ f)^*
\sigma_\phi$.  The map $s\mapsto\psi_s$ is a path of quasiconformal
homeomorphisms all of which fix the normalized triple, and at each time the
marking of $A$ is the identity marking; by
Lemma~\ref{lem:normalization} applied with the marking data fixed, the
endpoints $[\psi_0]_A$ and $[\psi_1]_A$ agree in $\Teich(\Sph,A)$.
\end{proof}

\begin{lemma}[holomorphy; see \cite{Hubbard}, and the discussion in
\cite{EpsteinTransversality}, \S2]\label{lem:holomorphic}
$f^{*}:\Teich(\Sph,B)\to\Teich(\Sph,A)$ is holomorphic.  In Beltrami
coordinates (Beltrami differentials vanishing near the marked points,
normalized by three marked points) the map is given by the linear--algebraic
pullback formula $(f^{*}\mu)(z)=\mu(f(z))\,
\overline{f'(z)}/f'(z)$ away from critical points and $0$ at critical
points, which depends complex--affinely on $\mu$.
\end{lemma}

\begin{remark}
Holomorphy will only be used to guarantee that $\D_f$ is a complex
analytic subvariety of $\Teich(\Sph,B)$ and that the image of each of its
components under $\pi:\Teich(\Sph,B)\to\M_B$ is an analytic subset; the
reader willing to take this as known may skip the proof, which is standard.
\end{remark}

\subsection{The dynamical characterisation of \texorpdfstring{$\D_f$}{Df}}

Let $\mu_0$ denote the zero Beltrami differential, and for a
quasiconformal $\phi$ write $\mu_\phi:=\phi^{*}\mu_0$.

\begin{proposition}[Dynamical characterisation]\label{prop:dynchar}
Let $\tau=[\phi]_B\in\Teich(\Sph,B)$, and let $s=\pi(\tau)\in\M_B$ be the
associated ordered configuration.  Then $\tau\in\D_f$ if and only if there
exists a quasiconformal homeomorphism $u:\Sph\to\Sph$ with $u|_A=\id|_A$
and $u\sim_A\id$, such that
\begin{equation}\label{eq:dynchar}
g:=\phi\circ f\circ u\circ\phi^{-1}:\widehat\C\to\widehat\C
\end{equation}
is a rational map.  In that case $g$ is the unique quadratic rational map
with superattracting $4$--cycle carrying the configuration $\phi(A)$ to
itself in the cyclic order, and with second critical value
$\phi\circ f\circ u(c_2)$; in particular if $\pi(\tau)$ is the base
configuration then $g=f$ and
\begin{equation}\label{eq:centralizer}
\phi\circ f\circ u=f\circ\phi .
\end{equation}
\end{proposition}

\begin{proof}
($\Leftarrow$) Suppose $u\sim_A\id$ and $g=\phi f u\phi^{-1}$ is rational.
Set $\Phi:=\phi\circ u^{-1}$, so that
$\mu_\Phi=(u^{-1})^{*}\mu_\phi$ and
$$g=\phi f \Phi^{-1},\qquad
g^{*}\mu_0=(\Phi^{-1})^{*}\, f^{*}\,\phi^{*}\mu_0=(\Phi^{-1})^{*}f^{*}\mu_\phi .$$
Since $g$ is rational, $g^{*}\mu_0=0$, and the Beltrami coefficient of
$\phi\circ f$ equals $f^{*}\mu_\phi=(\phi f)^{*}\mu_0$; the vanishing of
$(\Phi^{-1})^{*}(f^{*}\mu_\phi)$ means that the solution $\psi$ of
$\psi^{*}\mu_0=f^{*}\mu_\phi$ differs from $\Phi$ by a conformal map:
$\Phi=M\circ\psi$ for some M\"obius $M$.  Then
$\phi = \Phi\circ u = M\psi u$, which, by
Definition~\ref{def:teich}, means $[\phi]_A=[\psi]_A=f^{*}[\phi]_B$; i.e.\
$i_*\tau=f^{*}\tau$ and $\tau\in\D_f$.

($\Rightarrow$) Suppose $f^{*}\tau=i_*\tau$.  By
Definition~\ref{def:pullback} and Definition~\ref{def:teich} there are a
M\"obius $M$ and a quasiconformal $w\sim_A\id$ with $\phi=M\circ\psi\circ w$,
where $\psi^{*}\mu_0=f^{*}\mu_\phi$.  Put $u:=w^{-1}$ and
$\Phi:=\phi\, u^{-1}=M\psi$; then $\mu_\Phi=f^{*}\mu_\phi$, and the map
$g:=\phi f\Phi^{-1}$ satisfies
$g^{*}\mu_0=(\Phi^{-1})^{*}f^{*}\mu_\phi=(\Phi^{-1})^{*}\mu_\Phi
=(\Phi\Phi^{-1})^{*}\mu_0=\mu_0$, so $g$ is rational.  Since
$u|_A=\id|_A$, the map $g$ carries $\phi(a)$ to $\phi(f(a))$ for each
$a\in A$, i.e.\ $g$ maps the configuration $\phi(A)$ through the marked
cycle.  Uniqueness of $g$ in the normalized family is
Lemma~\ref{lem:family}(\ref{item:unique}) below.  If moreover
$\pi(\tau)$ is the base configuration then after a M\"obius normalization
we may take $\phi|_B=\id|_B$, so $g$ has the base cycle and base critical
value configuration; uniqueness forces $g=f$ and \eqref{eq:centralizer}.
\end{proof}

\subsection{The image of \texorpdfstring{$\D_f$}{Df} in moduli}

Recall from the problem statement that the image of $\D_f$ in moduli is
the set of conjugacy classes of quadratic rational maps with the properties
described above.  For completeness we indicate the proof; it also fixes
notation used later.

\begin{proposition}\label{prop:image}
Let $\pi:\Teich(\Sph,B)\to\M_B$ be the quotient.  Then:
\begin{enumerate}
\item $\pi(\D_f)$ is a complex analytic subset of $\M_B$, contained in the
\emph{stratum} $\S$ of configurations arising from quadratic rational maps
with a superattracting $4$--cycle marked in order, together with the second
critical value marked;
\item the base point $\tau_0=[\id]_B$ lies in $\D_f$;
\item the stratum $\S$ is an irreducible complex curve;
\item $\pi(\D_f)=\S$.
\end{enumerate}
\end{proposition}

\begin{proof}
(1) A point of $\pi(\D_f)$ is, by Proposition~\ref{prop:dynchar}, the
configuration of a quadratic rational map with a superattracting $4$--cycle
and the second critical value marked, i.e.\ a point of $\S$; and
$\pi(\D_f)$ is analytic because locally $\pi$ is biholomorphic and
$\D_f$ is analytic (Lemma~\ref{lem:holomorphic}).

(2) $\tau_0=[\id]_B$: pulling back the standard structure by the rational
map $f$ returns the standard structure, and the identity marking witnesses
$f^{*}\tau_0=i_*\tau_0$.

(3) This is Corollary~\ref{cor:stratumcurve}.

(4) By Proposition~\ref{prop:dynchar}(2) and
Lemma~\ref{lem:family}, the stratum is an irreducible curve; an analytic
subset of an irreducible curve is either finite or the whole curve minus a
finite set.  Since $\tau_0\in\D_f$, near $\tau_0$ the analytic set $\D_f$
has complex dimension at least $\dim\Teich(\Sph,B)-\dim\Teich(\Sph,A)=1$
(the fibre of the holomorphic map $f^{*}-i_{*}$ over the diagonal), so
$\pi(\D_f)$ contains a nonempty open piece of $\S$, hence is cofinite in
$\S$; being also analytic (hence closed in $\S$ after including
accumulation points), it equals $\S$.
\end{proof}

\section{Dehn twists through the double cover, and the preservation
criterion}\label{s:twists}

Throughout this section, $\gamma$ is an essential simple closed curve in
$(\Sph,B)$, with sides $D_X,D_Y$ and corresponding subsets
$X,Y\subseteq B$ ($X\sqcup Y=B$, $|X|,|Y|\ge2$).  Since
$V(f)=\{p_1,v\}\subseteq B$ and $f$ has no other critical values, the map
$$f:\Sph\setminus f^{-1}(B)\longrightarrow\Sph\setminus B$$
is an (unbranched) covering of degree $2$.

\subsection{The preimage multicurve and its degrees}

\begin{lemma}\label{lem:preimage}
Let $N\cong\overline{\mathbb{D}}\setminus\{0\}$ be a closed annular
neighbourhood of $\gamma$ in $\Sph\setminus B$.  Each component
$\tilde N_i$ of $f^{-1}(N)$ is an annulus mapping onto $N$ as a covering of
some degree $d_i\in\{1,2\}$, and the multicurve
$f^{-1}(\gamma)=\bigcup_i\tilde\gamma_i$
has components $\tilde\gamma_i=\tilde N_i\cap f^{-1}(\gamma)$, each an
essential curve in $(\Sph,f^{-1}B)$, with
$\deg(\tilde\gamma_i\to\gamma)=d_i$ and $\sum_i d_i=2$.
Moreover:
\begin{enumerate}
\item[(i)] if the two branch values $p_1,v$ lie on the \emph{same} side of
$\gamma$, then $f^{-1}(\gamma)$ has $d=1+1$: two components, and
$f^{-1}(D_X)$ (resp.\ $f^{-1}(D_Y)$) is the union of the two disks over the
side free of branch values, while over the side containing both branch
values the preimage is an annulus;
\item[(ii)] if $p_1$ and $v$ lie on \emph{opposite} sides, then
$f^{-1}(\gamma)$ has a single component of degree $2$, and each side's
preimage is a disk, branched over the one branch value it contains.
\end{enumerate}
\end{lemma}

\begin{proof}
The monodromy of the double cover around $\gamma$ is trivial in case (i)
(a loop enclosing both or neither branch value has trivial monodromy) and a
transposition in case (ii).  The classification of coverings of the annulus
$N$ gives the stated annuli $\tilde N_i$ and degrees.  Each
$\tilde\gamma_i$ is the core of $\tilde N_i$, hence essential in
$\Sph\setminus f^{-1}(B)$.  The description of the preimages of the sides
follows from the Riemann--Hurwitz and monodromy computations: a double
cover of a disk branched at one interior point is a disk, and branched at
two interior points is an annulus.
\end{proof}

\subsection{Lifting powers of the twist}

\begin{lemma}[Lift of a power of a twist]\label{lem:lift}
Keep the notation of Lemma~\ref{lem:preimage}, and let $k\ge1$.
\begin{enumerate}
\item The power $T_\gamma^{k}$ lifts through the cover
$f:(\Sph,f^{-1}B)\to(\Sph,B)$---that is, there exists a homeomorphism
$\widehat{T}$ of $(\Sph,f^{-1}B)$ fixing $f^{-1}(B)$ pointwise with
$f\circ\widehat{T}=T_\gamma^{k}\circ f$---if and only if $d_i\mid k$ for
every $i$.
\item When the lift exists, it is unique and equals the multitwist
\begin{equation}\label{eq:liftformula}
\widehat{T_\gamma^{k}}\ =\ \prod_i T_{\tilde\gamma_i}^{\,k/d_i}
\qquad\text{in }\Mod(\Sph,f^{-1}B).
\end{equation}
\end{enumerate}
\end{lemma}

\begin{proof}
Model $N=S^1\times[0,1]$ with coordinates $(s,t)$ and, on a component
$\tilde N_i$ mapping with degree $d_i$, write
$f(s,t)=(d_i\,s,t)$.  Let $T_\gamma$ be represented by the twist supported
in the interior of $N$, with profile $\theta:[0,1]\to[0,1]$,
$\theta(0)=\theta(1)\equiv0\pmod{\Z}$; in the usual normalization the
profile increases monotonically from $0$ to $1$ (the ``full twist'',
which is the identity on $\partial N$ because a full rotation of the
circle is the identity on the nose).  A lift $\widehat{T}$ must satisfy
$f(\widehat{T}(s,t))=(d_is+\theta(t),t)$, so
$\widehat T(s,t)=(s+\theta(t)/d_i,t)$ on $\tilde N_i$.  For this to extend
to a homeomorphism of the \emph{closed} annulus $\tilde N_i$ which is the
identity on $\partial\tilde N_i$ we need $\theta(t)/d_i\in\Z$ at $t=1$,
i.e.\ $1/d_i\in\Z$, i.e.\ $d_i\mid1$; for $T_\gamma^{\,k}$ the profile is
$k\theta$, and the condition becomes $k/d_i\in\Z$.  This proves (1), and
also exhibits the lift, whose restriction to $\tilde N_i$ is a twist of
$\tilde N_i$ with profile $k\theta/d_i$, i.e.\ the Dehn twist
$T_{\tilde\gamma_i}^{k/d_i}$; the lift extends by the identity outside
$f^{-1}(N)$.  Uniqueness: two lifts differ by the deck involution of the
\emph{unbranched} cover, which interchanges the two points of a general
fibre of $f^{-1}(B)$; since the lift must fix $f^{-1}(B)$ pointwise, the
deck factor is trivial.
\end{proof}

\subsection{Intertwining of the pullback with the mapping class group}

\begin{lemma}\label{lem:intertwine2}
Let $\gamma$ be as above and $k\ge1$ with $d_i\mid k$ for all $i$.  Then
for every $\tau\in\Teich(\Sph,B)$,
\begin{equation}\label{eq:intertwine2}
f^{*}\bigl(T_\gamma^{\,k}\cdot\tau\bigr)
=\Theta_k(\gamma)\cdot f^{*}\tau,
\end{equation}
where $\Theta_k(\gamma)\in\Mod(\Sph,A)$ is the image of
$\widehat{T_\gamma^{\,k}}$ under the forgetful homomorphism
$\pi_1^{f^{-1}B,\,A}:\Mod(\Sph,f^{-1}B)\to\Mod(\Sph,A)$.
\end{lemma}

\begin{proof}
Write $\tau=[\phi]_B$ and let $h$ represent $T_\gamma^{\,k}$.  Then
$T_\gamma^{\,k}\cdot\tau=[\phi\circ h^{-1}]_B$, whose Beltrami coefficient
is $(h^{-1})^{*}\mu_\phi$.  Using the lift $\widehat h$ of $h$ and the
identity $h^{-1}\circ f=f\circ\widehat h^{-1}$ (from
$f\widehat h=hf$), we get
$$f^{*}\bigl((h^{-1})^{*}\mu_\phi\bigr)
=(h^{-1}f)^{*}\mu_\phi=(\widehat h^{-1})^{*}\,f^{*}\mu_\phi,$$
so the pulled--back structure is the $f^{*}\tau$--structure pre--composed
with $\widehat h^{-1}$.  Unwinding
Definition~\ref{def:pullback} (the marking of the pulled--back structure is
the identity marking of the underlying set), this means exactly
$f^{*}(T_\gamma^{\,k}\tau)=[\psi\circ\widehat h^{-1}]_A
=\pi_1^{f^{-1}B,A}(\widehat h)\cdot[\psi]_A$,
where $[\psi]_A=f^{*}\tau$.  Well--definedness against the choice of lift
is Lemma~\ref{lem:lift}(2), and against the representative of $\tau$ it
follows as in Lemma~\ref{lem:welldefined}.
\end{proof}

\subsection{The preservation criterion}

Recall the forgetful homomorphism
$\pi_1:=\pi_1^{B,A}:\Mod(\Sph,B)\to\Mod(\Sph,A)$ of \eqref{eq:forgethom}.
For a curve $\gamma$ in $(\Sph,B)$ write $\bar\gamma$ for the same curve
viewed in $(\Sph,A)$ (forgetting $v$); it is essential in $(\Sph,A)$ iff
both sides of $\gamma$ contain at least two points of $A$, in which case
$\pi_1(T_\gamma^{\,k})=T_{\bar\gamma}^{\,k}$, while if some side contains at
most one point of $A$ then $\pi_1(T_\gamma^{\,k})=\id$ by
Lemma~\ref{lem:peripheral}.

\begin{theorem}[Preservation Criterion]\label{thm:crit}
Let $\gamma$ be an essential curve in $(\Sph,B)$, with preimage components
$\tilde\gamma_i$ and degrees $d_i$ as in Lemma~\ref{lem:preimage}, and let
$k\ge1$ satisfy $d_i\mid k$ for all $i$.  Assume that the stabilizer in
$\Mod(\Sph,A)$ of the (open) set $i_*(\D_f)\subseteq\Teich(\Sph,A)$ is
trivial.  Then
$$T_\gamma^{\,k}\cdot\D_f=\D_f
\quad\Longleftrightarrow\quad
\Theta_k(\gamma)=\pi_1(T_\gamma^{\,k})\ \text{ in }\Mod(\Sph,A).$$
Here, by Lemmas \ref{lem:peripheral} and \ref{lem:lift},
\begin{equation}\label{eq:theta}
\Theta_k(\gamma)=\prod_{i\,:\,\tilde\gamma_i\text{ essential in }(\Sph,A)}
T_{\tilde\gamma_i}^{\,k/d_i},
\end{equation}
the product over those components of $f^{-1}(\gamma)$ that separate $A$
into two parts of size $\ge2$.
\end{theorem}

\begin{proof}
By Lemma~\ref{lem:intertwine1}, for all $\tau$,
$i_*(T_\gamma^{\,k}\tau)=\pi_1(T_\gamma^{\,k})\,i_*(\tau)$; by
Lemma~\ref{lem:intertwine2}, for all $\tau\in\D_f$,
$f^{*}(T_\gamma^{\,k}\tau)=\Theta_k(\gamma)\,f^{*}\tau
=\Theta_k(\gamma)\,i_*(\tau)$.
Hence for $\tau\in\D_f$:
$$T_\gamma^{\,k}\tau\in\D_f
\iff \Theta_k(\gamma)\,i_*(\tau)=\pi_1(T_\gamma^{\,k})\,i_*(\tau)
\iff \bigl(\Theta_k(\gamma)^{-1}\pi_1(T_\gamma^{\,k})\bigr)\text{ fixes
}i_*(\tau).$$
Since $i_*(\D_f)$ is a nonempty open subset of the $1$--dimensional space
$\Teich(\Sph,A)$ (Proposition~\ref{prop:image}(4) and
Lemma~\ref{lem:family}(\ref{item:surjective})), and $\Mod(\Sph,A)$ acts
freely there, the last condition holds for all $\tau\in\D_f$ iff
$\Theta_k(\gamma)=\pi_1(T_\gamma^{\,k})$.  The case of equality gives
$T_\gamma^{\,k}\D_f\subseteq\D_f$, and equality follows by applying the
same to $T_\gamma^{-\,k}$.
\end{proof}

\begin{remark}
The triviality of the stabilizer of $i_*(\D_f)$ is automatic in our
situation: $i_*(\D_f)$ contains an open subset of $\Teich(\Sph,A)$
(namely the set of cross--ratios of $4$--cycles realized by stratum maps,
cf.\ \S\ref{s:stratum}), and the pure mapping class group of the
$4$--punctured sphere acts freely on Teichm\"uller space.
\end{remark}

\section{The stratum: an explicit algebraic model}\label{s:stratum}

\subsection{The normalized family}

\begin{lemma}\label{lem:family}
Let $t\in\widehat\C\setminus\{0,1\}$.  There is a unique quadratic rational
map $f_t$ with $f_t(0)=1,\ f_t(1)=\infty,\ f_t(\infty)=t,\ f_t(t)=0$ and
with $0$ as a critical point; it is
\begin{equation}\label{eq:family}
f_t(z)=\frac{t\,z^{2}-(t^{2}+t+1)\,z+t^{2}+t}
{z^{2}-(t^{2}+t+1)\,z+t^{2}+t},
\end{equation}
and:
\begin{enumerate}
\item\label{item:unique} $f_t$ is the \emph{unique} quadratic rational map
$g$ with $g(0)=1,\ g(1)=\infty,\ g(\infty)=t,\ g(t)=0$ and a critical
point at $0$;
\item the second critical point and its value are
\begin{equation}\label{eq:c2v}
c_2(t)=\frac{2t(t+1)}{t^{2}+t+1},
\qquad
v(t)=f_t(c_2(t))
=\left(\frac{t^{2}-t-1}{t^{2}+t-1}\right)^{\!2};
\end{equation}
\item\label{item:surjective} $f_t$ has exact period $4$ at $0$ for
$t\in\C\setminus\{0,1\}$, and $t\mapsto f_t$ is a bijection from
$\C\setminus\{0,1\}$ onto quadratic rational maps with a marked critical
point of exact period $4$, normalized by the cycle positions
$0\mapsto1\mapsto\infty\mapsto t\mapsto0$;
\item the marked configuration of $f_t$ is
$(0,1,\infty,t,v(t))$ for $(c_1,p_1,p_2,p_3,v)$.
\end{enumerate}
\end{lemma}

\begin{proof}
Writing $g(z)=\frac{az^2+bz+c}{dz^2+ez+g_0}$ and imposing
$g(0)=1$ ($c=g_0$), $g(1)=\infty$ ($d+e+g_0=0$),
$g(\infty)=t$ ($a=td$), $g(t)=0$ ($at^2+bt+c=0$) and
$g'(0)=0$ ($bg_0=ce$, i.e.\ $b=e$), and normalizing $g_0=1$, one solves the
linear conditions to obtain \eqref{eq:family} (a direct substitution
verifies each condition).  Since the five conditions determine all
coefficients up to the common scalar, uniqueness (\ref{item:unique}) holds;
this is the uniqueness statement used in
Propositions~\ref{prop:dynchar} and \ref{prop:image}.  The critical points
of $f_t$ are the zeros of $N'D-ND'$, which factorizes as
$(t-1)\,z\,\bigl(2(t^2+t)-(t^2+t+1)\,z\bigr)$, giving $0$ and
$c_2(t)$ as in \eqref{eq:c2v}; substituting gives $v(t)$, using
$(t^2-t-1)^2=t^4-2t^3-t^2+2t+1$ and
$(t^2+t-1)^2=t^4+2t^3-t^2-2t+1$.
For (\ref{item:surjective}), the orbit $0\mapsto1\mapsto\infty\mapsto
t\mapsto0$ has exact period $4$ whenever the four points are distinct and
the earlier return values $f_t^{\circ j}(0)=1,\infty,t$ for $j=1,2,3$ are
different from $0$; this fails only at $t\in\{0,1\}$ (at $t=0,-1$ the
map degenerates; at $t=1$ it is constant).  Finally a M\"obius change of
coordinates taking a marked map to normal form produces some $t$, and the
marked point $p_3$ is that $t$.
\end{proof}

\begin{corollary}\label{cor:stratumcurve}
The stratum of ordered configurations
$(c_1,p_1,p_2,p_3,v)=(0,1,\infty,t,v)$ arising from quadratic rational
maps with a critical point of exact period $4$ and second critical value off
the cycle is the graph
$$\S=\{(t,v)\ :\ v=v(t)\}\ \subset\ \C^2\cong\M_B^{\mathrm{ord}},$$
an irreducible rational curve, punctured at the parameters $t$ where the
five--tuple $(0,1,\infty,t,v(t))$ fails to have distinct entries or the map
degenerates.
\end{corollary}

\begin{definition}\label{def:cusps}
The \emph{cusps} are the parameters $t^{*}$ at which $v(t)$ collides with a
point of the cycle or with $t$, or the cycle degenerates:
\begin{align*}
&v=0:\quad a_\pm=\tfrac{1\pm\sqrt5}{2} &&\text{(note }v(t)=0\iff
t^2-t-1=0);\qquad
v=1:\ t=-1;\\
&v=\infty:\quad b_\pm=\tfrac{-1\pm\sqrt5}{2};\qquad &&
v=t:\ \text{the roots of } t^{5}+t^{4}-3t^{3}-3t^{2}+3t+1;\\
&t\in\{0,1,-1,\infty\}:\ \text{cycle degeneration}.
\end{align*}
We call $a_\pm$ and $b_\pm$ the \emph{golden cusps}.  At a golden cusp
$t^{*}$, $v(t^{*})$ equals $c_1$ (for $a_\pm$) or $p_2$ (for $b_\pm$), the
second critical point $c_2(t^{*})$ lands on the cycle
(one checks $c_2(a_\pm)=a_\pm$, $c_2(b_\pm)=\frac{4}{3\mp\sqrt5}$ and
$f_{t^{*}}(c_2(t^{*}))=v(t^{*})$), and the map $f_{t^{*}}$ is
\emph{postcritically finite} with postcritical set exactly the cycle $A$.
\end{definition}

\begin{lemma}[Braid monodromy of the cusps]\label{lem:monodromytable}
Let $\ell_{t^{*}}$ be a small positively oriented loop in the $t$--plane
around a cusp $t^{*}$.  The braid monodromy $b(\ell_{t^{*}})\in
\Mod(\Sph,B)$ (the deck displacement of the lifted path of configurations
in $\Teich(\Sph,B)$, cf.\ Lemma~\ref{lem:normalization}) is:
\begin{center}
\begin{tabular}{c|c|c}
cusp $t^{*}$ & colliding pair & $b(\ell_{t^{*}})$\\ \hline
$0,\ 1,\ \infty$ & $\{p_3,c_1\},\ \{p_3,p_1\},\ \{p_3,p_2\}$
& $T_{\{c_1,p_3\}}^{\pm1},\ T_{\{p_1,p_3\}}^{\pm1},\ T_{\{p_2,p_3\}}^{\pm1}$\\
$-1$ & $\{v,p_1\}$ & $T_{\{p_1,v\}}^{\pm1}$\\
$a_\pm$ & $\{v,c_1\}$, twice & $T_{\{c_1,v\}}^{\pm2}$\\
$b_\pm$ & $\{v,p_2\}$, twice & $T_{\{p_2,v\}}^{\pm2}$\\
roots of $v(t)=t$ & $\{v,p_3\}$ & $T_{\{p_3,v\}}^{\pm1}$
\end{tabular}
\end{center}
In particular, for the golden cusps the factor $2$ in the exponent is
forced by \eqref{eq:c2v}: near $t^{*}=a_\pm$ one has
$v(t)=C\,(t-t^{*})^{2}+O\bigl((t-t^{*})^3\bigr)$ with $C\ne0$, so one turn
of $t$ around $t^{*}$ produces two turns of $v$ around $c_1$; similarly at
$b_\pm$, where $v$ has a double pole.
\end{lemma}

\begin{proof}
The configuration $(0,1,\infty,t,v(t))$ moves with $t$; only the two
strands that collide at the cusp braid nontrivially, and the winding number
of the moving strand around the marked point it approaches is the local
degree of $v$ (resp.\ of $t$) at the cusp.  At $a_\pm$ the numerator
$t^2-t-1$ of $v$ has a simple zero, so $v$ vanishes to order exactly $2$;
at $b_\pm$ the denominator $t^2+t-1$ has a simple zero, so $v$ has a pole
of order $2$.  The remaining entries have order $1$.
\end{proof}

\begin{remark}[The four PCF limits and their invariant equator]
\label{rem:carousal}
Each golden cusp $t^{*}\in\{a_\pm,b_\pm\}$ is real, and $f_{t^{*}}$ is a
\emph{real} rational map (real coefficients) whose two critical points
$0$ and $c_2(t^{*})$ are real and lie on the circle
$\R\cup\{\infty\}$.  For these four maps the circle
$J:=\R\cup\{\infty\}$ is a fully invariant Jordan curve: since both
critical values $p_1,v(t^{*})$ lie on $J$, the two ``laps'' of
$f_{t^{*}}|_J$ each map onto $J$, so
$f_{t^{*}}^{-1}(J)=J$, and $f_{t^{*}}$ maps each component of
$\widehat\C\setminus J$ conformally onto a component of
$\widehat\C\setminus J$.  Thus the four golden maps admit an invariant
Jordan curve through the postcritical set (a ``Levy curve'' or
\emph{carousal}).  This classical structure is the geometric backbone of
the class--monodromy analysis of \S\ref{s:components}.
\end{remark}

\section{The twist \texorpdfstring{$T_\gamma^{\,2}$}{T2} preserves
\texorpdfstring{$\D_f$}{Df}}\label{s:preserve}

Here is the curve promised in the Main Theorem.

\begin{definition}\label{def:gamma}
Let $\gamma$ be a simple closed curve in $(\Sph,B)$ bounding a disk $D$
with $D\cap B=\{c_1,v\}$; thus the complementary disk contains
$\{p_1,p_2,p_3\}$.  Such curves exist (join $c_1$ to $v$ by an arc
avoiding $B\setminus\{c_1,v\}$ and take a regular neighbourhood of the
arc) and any two are isotopic in $(\Sph,B)$, since they induce the same
partition of $B$.
\end{definition}

\begin{theorem}\label{thm:preserve}
$T_\gamma^{\,2}\cdot\D_f=\D_f$.
\end{theorem}

\begin{proof}
We verify the hypotheses and the equality of
Theorem~\ref{thm:crit} with $k=2$.

\emph{Step 1: the preimage of $\gamma$.}  The branch values of $f$ are
$p_1$ and $v$; of these, $v\in D$ and $p_1\notin\overline D$, so
Lemma~\ref{lem:preimage}(ii) applies: $f^{-1}(\gamma)$ is a single curve
$\tilde\gamma$ mapping onto $\gamma$ with degree $d=2$, and
$f^{-1}(D)$ is a disk (the double cover of $D$ branched over the single
branch value $v\in D$).  Concretely
$$f^{-1}(D)\ \supset\ f^{-1}\bigl(D\cap B\bigr)
=f^{-1}(\{c_1,v\})=\{p_3,w_0,c_2\},$$
where $w_0$ denotes the second preimage of $c_1$ (which is not a point of
$B$), while
$$f^{-1}\bigl(B\setminus\{c_1,v\}\bigr)
=f^{-1}(\{p_1,p_2,p_3\})
=\{c_1,w_1\}\cup\{p_1,w_2\}\cup\{p_2,w_3\}$$
lies in the complementary disk.  Hence $\tilde\gamma$ separates
$\{p_3,w_0,c_2\}$ from the remaining six points of $f^{-1}(B)$.

\emph{Step 2: $\tilde\gamma$ is peripheral in $(\Sph,A)$.}  Indeed
$\tilde\gamma$ bounds a disk meeting $A$ in the single point $p_3$
(the other two enclosed points $w_0,c_2$ do not belong to $A$).  By
Lemma~\ref{lem:peripheral}, the class of $T_{\tilde\gamma}$ in
$\Mod(\Sph,A)$ is trivial; hence in \eqref{eq:theta} the single factor is
trivial and
$$\Theta_2(\gamma)=\id\in\Mod(\Sph,A).$$

\emph{Step 3: $\bar\gamma$ is peripheral in $(\Sph,A)$.}  The curve
$\gamma$, with the point $v$ forgotten, bounds a disk meeting
$A=\{c_1,p_1,p_2,p_3\}$ only in $c_1$; by
Lemma~\ref{lem:peripheral},
$$\pi_1(T_\gamma^{\,2})=\id\in\Mod(\Sph,A).$$

\emph{Step 4: conclusion.}  Since $d=2$ divides $k=2$,
Lemma~\ref{lem:lift} provides the lift
$\widehat{T_\gamma^{\,2}}=T_{\tilde\gamma}^{\,2/2}=T_{\tilde\gamma}$ in
$\Mod(\Sph,f^{-1}B)$, and
Lemmas \ref{lem:intertwine1}, \ref{lem:intertwine2} apply.  Since
$\Theta_2(\gamma)=\id=\pi_1(T_\gamma^{\,2})$, the criterion of
Theorem~\ref{thm:crit} gives $T_\gamma^{\,2}\D_f=\D_f$.

For the variants: if $\gamma$ encloses exactly $\{p_2,v\}$, then
$f^{-1}(\gamma)$ is again a single degree--$2$ curve enclosing
$\{p_1,w_2,c_2\}$, peripheral in $(\Sph,A)$ around $p_1$, and
$\bar\gamma$ is peripheral around $p_2$; if $\gamma$ encloses exactly
$\{p_3,v\}$, the preimage is a degree--$2$ curve enclosing
$\{p_2,w_3,c_2\}$, peripheral around $p_2$, and $\bar\gamma$ is peripheral
around $p_3$.  The same computation applies verbatim.
\end{proof}

\begin{remark}[why other curves fail]
It is instructive to record why nearby--looking curves do \emph{not} work,
as this justifies the choice of $\gamma$:
\begin{itemize}
\item a curve enclosing $\{c_1,p_1\}$ has a connected degree--$2$
preimage enclosing $\{c_1,p_3\}$, which is \emph{essential} in
$(\Sph,A)$ but with the wrong partition; the exponents $k/2$ and $k$ can
never match;
\item a curve enclosing $\{p_1,p_3\}$ (the ``alternating'' pair) has a
connected degree--$2$ preimage enclosing $\{c_1,p_2\}$ with the
\emph{right} partition, but the lift of $T_\gamma^{\,k}$ is
$T_{\tilde\gamma}^{\,k/2}$, which can never equal $T_{\bar\gamma}^{\,k}$
for $k\neq0$;
\item a curve enclosing $\{p_1,v\}$ has a preimage which is an annulus
whose two boundary curves are, up to isotopy in $(\Sph,A)$, twists about
curves enclosing pairs of cycle points, none of which is peripheral; the
two sides of the criterion never agree.
\end{itemize}
Thus among the ``pair--twists'' exactly the three enclosing
$\{x,v\}$ with $x\in\{c_1,p_2,p_3\}$ have a chance, and they all work
after squaring.
\end{remark}

\section{Components of \texorpdfstring{$\D_f$}{Df} and the infinite orbit}
\label{s:components}

Fix the base data: after a M\"obius conjugacy we assume
$f=f_{t_0}$ in the family \eqref{eq:family} for some
$t_0\in\C\setminus(\P\cup\{1\})$ with $v(t_0)\notin\{0,1,\infty,t_0\}$, so
that the marked configuration is $B=(0,1,\infty,t_0,v_0)$ with
$v_0=v(t_0)$.  Let
$$\tau_0:=[\id]_B\in\D_f$$
be the base point, and let
$$\sigma_1:=T_{\{c_1,v\}}^{2},\qquad
\sigma_2:=T_{\{p_2,v\}}^{2},\qquad
\sigma_3:=T_{\{p_3,v\}}^{2}\ \in\Mod(\Sph,B),$$
so that by Theorem~\ref{thm:preserve}, $\sigma_i\cdot\D_f=\D_f$; note that
$\sigma_1=T_\gamma^{\,2}$ for our $\gamma$.

\subsection{The fibre of \texorpdfstring{$\D_f\to\S$}{Df→S} over the base}

\begin{definition}\label{def:Z}
Let $Z\le\Mod(\Sph,B)$ be the set of classes $[\phi]$ with
$\phi|_B=\id_B$ that admit a quasiconformal $u\sim_A\id$, $u|_A=\id$,
with $\phi\circ f\circ u=f\circ\phi$; i.e.\ the group of marked
self--conjugacy classes of $f$ ``up to $A$--isotopy''.
\end{definition}

\begin{proposition}\label{prop:fibre}
\begin{enumerate}
\item $Z$ is a subgroup of $\Mod(\Sph,B)$, and
$z\mapsto z\cdot\tau_0$ identifies $Z$ with the fibre
$F_0:=\D_f\cap\pi^{-1}(s_0)$ of $\pi:\D_f\to\S$ over the base
configuration $s_0$.
\item For $[\phi]\in Z$ the lift $\widehat\phi$ of $\phi$ through
$f$ fixes $f^{-1}(B)$ pointwise, and
$$Z=\bigl\{[\phi]\in\Mod(\Sph,B)\ :\ \exists\text{ lift }\widehat\phi
\text{ of }\phi\text{ with }\widehat\phi\sim_A\phi\bigr\}.$$
In particular $Z$ coincides with the subgroup on which the two
homomorphisms $\pi_1$ and $\Theta:=\pi_1^{f^{-1}B,\,A}\circ\widehat{(\cdot)}$
of \S\ref{s:twists} agree.
\end{enumerate}
\end{proposition}

\begin{proof}
(1) If $\phi_1 f u_1=f\phi_1$ and $\phi_2fu_2=f\phi_2$ then
$\phi_1\phi_2 f(u_2\circ\phi_2^{-1}u_1\phi_2)=f\phi_1\phi_2$ with
$u_2':=\phi_2^{-1}u_1\phi_2\sim_A\id$, so $Z$ is closed under composition;
inverses are analogous.  By Lemma~\ref{lem:normalization} the fibre
$\pi^{-1}(s_0)$ consists of the classes $[\phi]_B$ with
$\phi|_B=\id_B$, and by Proposition~\ref{prop:dynchar} together with the
uniqueness in Lemma~\ref{lem:family}(\ref{item:unique}),
$[\phi]_B\in\D_f$ iff $\phi$ satisfies the defining equation of $Z$.

(2) From $\phi fu=f\phi$ we get $f\circ(\phi u^{-1})=\phi\circ f$, so
$\widehat\phi:=\phi u^{-1}$ is a lift of $\phi$; it fixes each $a\in A$
(since $u|_A=\id$ and $\phi|_B=\id$), and for $x\in f^{-1}(B)$ we have
$f(\widehat\phi(x))=\phi f(x)=f(x)$, so $\widehat\phi$ permutes each fibre
$f^{-1}(f(x))$; since $\widehat\phi$ already fixes the point of
$A\cap f^{-1}(f(x))$, it fixes both points of the fibre.  As
$u\sim_A\id$, $\widehat\phi=\phi u^{-1}\sim_A\phi$.  The converse is
immediate, and the last assertion is a restatement (this also recovers
Theorem~\ref{thm:crit} on $Z$).
\end{proof}

\subsection{Monodromy and the coset description of components}

Let $M\subseteq Z$ be the set of elements $g_\ell\in Z$ obtained as
follows: $\ell$ ranges over loops in the stratum $\S$ based at $s_0$ such
that the lift of $\ell$ to $\Teich(\Sph,B)$ starting at $\tau_0$ remains
in $\D_f$ for all time, and $g_\ell\in\Mod(\Sph,B)$ is the deck
transformation taking the initial point $\tau_0$ of the lifted path to its
endpoint.  (The endpoint lies in $F_0\cong Z$, so $g_\ell\in Z$.)

\begin{proposition}\label{prop:cosets}
\begin{enumerate}
\item $M$ is a subgroup of $Z$, and the twist $\sigma_i$ acts on
$\pi_0(\D_f)$; under the identification $F_0\cong Z$, the action of
$\sigma_1=T_\gamma^{\,2}$ on the fibre is by left multiplication, and
$\pi_0(\D_f)$ is identified with the set of right cosets $Z/M$ in such a
way that $\sigma_1$ acts by $zM\mapsto \sigma_1\,zM$.
\item The orbit of the coset $zM$ under $\langle\sigma_1\rangle$ is
infinite iff for every $k\ne0$,
\begin{equation}\label{eq:separation}
z^{-1}\sigma_1^{\,k}z\notin M .
\end{equation}
\end{enumerate}
\end{proposition}

\begin{proof}
(1) Concatenation of loops shows $M$ is a subgroup; the identification of
the fibre with $Z$ is Proposition~\ref{prop:fibre}; the deck displacement
of the lift of $\ell$ started at $z\tau_0$ is $z\,g_\ell$ because $z$,
belonging to $Z$, carries the $\tau_0$--lift to the $z\tau_0$--lift and
preserves $\D_f$; hence the monodromy acts on fibre--coordinates by right
multiplication, and two fibre points lie in the same component iff they
differ by an element of $M$.  The twist $\sigma_1\in Z$ maps components to
components, and on fibre--coordinates it acts by left multiplication
$z\mapsto\sigma_1z$.

(2) $\sigma_1^{\,k}zM=zM\iff z^{-1}\sigma_1^{-k}z\in M$; replacing $k$ by
$-k$ and noting $M$ is a group gives \eqref{eq:separation}.  Finitely many
cosets in the orbit would give, for some $k\ne0$ and some iterate,
$\sigma_1^{\,k}zM=zM$ for all $z$ in a fixed coset representative set;
by (1) this is exactly the negation of \eqref{eq:separation} for that $z$.
\end{proof}

\subsection{A free subgroup inside \texorpdfstring{$Z$}{Z}}

\begin{proposition}\label{prop:free}
The subgroup $F:=\langle\sigma_1,\sigma_2,\sigma_3\rangle\le Z$ is free of
rank $3$ with basis $\{\sigma_1,\sigma_2,\sigma_3\}$.
\end{proposition}

\begin{proof}
By Theorem~\ref{thm:preserve} (and its variants) each $\sigma_i$ lies in
$Z$ (equivalently: run the criterion of Theorem~\ref{thm:crit} for the
curves enclosing $\{c_1,v\},\{p_2,v\},\{p_3,v\}$ with $k=2$).  Let
$\Pi\cong\pi_1(\Sph\setminus A)\cong F_3$ be the point--pushing subgroup
of $\Mod(\Sph,B)$ (Lemma~\ref{lem:birman}, with the pushed point $v$); it
is freely generated by the pushes along small loops
$\beta_{c_1},\beta_{p_2},\beta_{p_3}$ around $c_1,p_2,p_3$, and by
Lemma~\ref{lem:birman} $\mathrm{push}(\beta_x)=T_{\{x,v\}}$, so
$\sigma_1=\mathrm{push}(\beta_{c_1}^{\,2})$,
$\sigma_2=\mathrm{push}(\beta_{p_2}^{\,2})$,
$\sigma_3=\mathrm{push}(\beta_{p_3}^{\,2})$.  By
Corollary~\ref{cor:pushfree} these three squares form a free basis of the
subgroup they generate, so $F\cong F_3$.
\end{proof}

\subsection{Class--monodromy at the golden cusps (dynamical input)}

For $t\in\S$ let us say that $f_t$ has the \emph{same combinatorial class
(rel.\ $A$) as $f$} if there exists a homeomorphism $h:(\Sph,B(t))\to
(\Sph,B)$ matching the marked data such that $h\circ f_t\circ
u\simeq_A f\circ h$ for some $u\sim_A\id$; by
Proposition~\ref{prop:dynchar}, along a lifted path
$\tau_s\in\D_f$ the maps $f_{\ell(s)}$ all have the same class as $f$,
and conversely.

The following is the dynamical core of the theorem.  Its content is that
one \emph{cannot} loop around a golden cusp while staying inside one
combinatorial class: the class changes by exactly the square twist
about the vanishing pair.

\begin{proposition}[Class--monodromy at the golden cusps]
\label{prop:input}
Let $t^{*}\in\{a_+,a_-\}$ (the case $\{b_+,b_-\}$ being identical with
$\sigma_1$ replaced by $\sigma_2$) and let $\ell$ be a loop in $\S$ based
at $t_0$ whose winding number around $t^{*}$ is nonzero.  Then the
combinatorial class of $f_{\ell(s)}$ (rel.\ $A$) is \emph{not} constant
along $\ell$.  Equivalently: the lift of $\ell$ to $\Teich(\Sph,B)$
starting at $\tau_0$ does not remain in $\D_f$; in particular, by
Proposition~\ref{prop:cosets} and Lemma~\ref{lem:monodromytable},
\begin{equation}\label{eq:nosigma}
\sigma_1^{\,k}\notin M\qquad\text{for all }k\ne0,
\end{equation}
and the analogous statement holds for $\sigma_2^{\,k}$.
\end{proposition}

\begin{proof}[Proof of Proposition~\ref{prop:input}]
We give the proof for $t^{*}=a_+$; the other cusps are treated by the
same argument with the evident changes of labels.  Recall
(Lemma~\ref{lem:family}, Definition~\ref{def:cusps}) that
$t^{*}=a_+=\tfrac{1+\sqrt5}{2}$ is real, that
$v(t)=\bigl(\frac{t^2-t-1}{t^2+t-1}\bigr)^2$ has a zero of exact order
$2$ at $t^{*}$, that $c_2(t^{*})=t^{*}=p_3(t^{*})$, and that
$F:=f_{t^{*}}$ is postcritically finite with postcritical set $A$.

\emph{Step 1 (invariant curve at the limit).}  $F$ is a real rational map
whose critical points $0$ and $t^{*}$ are real, so the circle
$J=\R\cup\{\infty\}$ contains both critical values $1$ and $0=v(t^{*})$;
consequently $F|_J:J\to J$ is a two--to--one covering (its two laps, cut at
the critical points, each map onto $J$), hence $F^{-1}(J)=J$ and $F$ maps
each component of $\widehat\C\setminus J$ conformally onto a component of
$\widehat\C\setminus J$ (Remark~\ref{rem:carousal}).  Thus $(F;A)$ carries
an invariant Jordan curve, with $A\subset J$ in the cyclic order
$0<1<t^{*}<\infty$.

\emph{Step 2 (the itinerary invariant).}  For a map $g=f_t$ whose
configuration is close to that of $F$, define its \emph{itinerary} as
follows.  By Step 1 and the structural stability of the invariant curve
under quasiconformal deformation (the curve $J$, together with its two
sides, moves by a quasiconformal isotopy of the sphere carrying marked
points to marked points; this is the standard ``$\lambda$--lemma'' transport
of the invariant curve along the stratum, cf.\ \cite{MSS} and, for
critically preperiodic maps specifically, \cite[Ch.~5]{BH} and
\cite{PilgrimTan}), for $t$ in a small simply connected neighbourhood
$U^{*}$ of $t^{*}$ with $t^{*}$ removed and cut along a wall, there is a
canonically associated Jordan curve $J_t$ through the four points of
$A(t)$, varying continuously with $t$, such that $f_t$ maps a
neighbourhood of the ``carousal'' structure consistently with $F$.
The \emph{itinerary} of $f_t$ is the symbol sequence recording, for the
orbit of the free critical value,
$$v(t),\quad f_t(v(t)),\quad f_t^{\circ 2}(v(t)),\ \ldots,$$
which of the two components of $\widehat\C\setminus J_t$ (and which of the
arcs of $J_t$ cut by $A(t)$) each of the first $N$ points occupies,
together with the winding datum of the arc from $v(t)$ to $f_t(v(t))$
around the marked points.  Two maps in the same combinatorial class
rel.\ $A$ have equal itineraries: a conjugacy $h$ with
$h f_t u\simeq_A f_{t'} h$ transports $J_t$ to $J_{t'}$ (both curves are
characterized intrinsically as the invariant curve of the carousal
structure, which is a class invariant), hence transports orbit to orbit
preserving the side/arc labels.

\emph{Step 3 (the itinerary changes around the cusp).}  Consider the two
real intervals adjacent to $t^{*}$: for
$t\in(t^{*},t^{*}+\varepsilon)$ one computes from \eqref{eq:c2v} that
$c_2(t)<t$ (e.g.\ $c_2(2)=\tfrac{12}7<2$), while for
$t\in(t^{*}-\varepsilon,t^{*})$ one has $c_2(t)>t$.  Consequently the
free critical point lies, relative to the cyclically ordered quadruple
$(0,1,t,\infty)$ on $J_t$, in the arc $(t,\infty)$ in the first case and in
the arc $(1,t)$ in the second; since $f_t(c_2(t))=v(t)$ and the map on the
carousal is fixed combinatorially, the itineraries of $f_t$ on the two
sides of $t^{*}$ (along the real axis) differ in their first symbol.
Hence the real interval $(t^{*},t^{*}+\varepsilon)$ and
$(t^{*}-\varepsilon,t^{*})$ carry \emph{different} classes, and the two
rays of the real axis are walls issuing from $t^{*}$.
Since the walls issuing from $t^{*}$ separate any small punctured
neighbourhood of $t^{*}$ into sectors on which the class is constant and
adjacent sectors carry distinct classes, every loop with nonzero winding
around $t^{*}$ crosses at least one wall and therefore exits the class.

\emph{Step 4 (from walls to $M$).}  By Step 3 and the discussion
preceding the proposition, a loop whose lift stays in $\D_f$ has zero
winding around each golden cusp.  Suppose now $g_\ell=\sigma_1^{\,k}\in M$
with $k\ne0$, where $\ell$ is the corresponding lifting loop.  The braid
$g_\ell$ is a pure $v$--pushing braid (it is a product of the three
generators $\sigma_1,\sigma_2,\sigma_3$ of $F$), so the $t$--strand is
unbraided; on the other hand, from Lemma~\ref{lem:monodromytable} and the
additivity of winding numbers, the total winding of $v$ around $c_1$
along $\ell$ equals $2k$, which is nonzero exactly when $k\ne0$; such
winding is contributed only by loops with nonzero winding around the
golden cusps $a_\pm$, contradicting Step 3.  This proves
\eqref{eq:nosigma}.
\end{proof}

\begin{remark}
Step 3 is the standard ``wake'' picture of critically marked strata: walls
(separatrices) issue from each cusp, and the combinatorial class is
constant precisely on the complement of the walls.  The explicit sign
change of $c_2(t)-t$ across $t^{*}=a_\pm$, together with the invariance of
the carousal $J$ for the four PCF limits (Step 1), makes the picture
completely concrete in our family; for the general theory of invariant
curves (``Levy curves'') and the rigidity of critically preperiodic maps we
refer to \cite{PilgrimTan, Kiwi, Rees}.  Note that the factor $2$ in
$v(t)=C(t-t^{*})^{2}+\cdots$ is consistent with the square twist
$\sigma_1=T_{\{c_1,v\}}^2$ of Lemma~\ref{lem:monodromytable}: a single
turn of the parameter is needed to produce the two turns of $v$ around
$c_1$, which is the unit of twisting measured by $\sigma_1$.
\end{remark}

\subsection{Proof of the infinite orbit}

\begin{theorem}\label{thm:infinite}
Let $z:=\sigma_2\in Z$.  Then for every $k\ne0$,
$z^{-1}\sigma_1^{\,k}z\notin M$.  Consequently, by
Proposition~\ref{prop:cosets}(2), the orbit of the component of
$\sigma_2\tau_0$ under the action of $T_\gamma^{\,2}=\sigma_1$ on
$\pi_0(\D_f)$ is infinite.  In particular $\D_f$ has infinitely many
connected components.
\end{theorem}

\begin{proof}
Suppose $z^{-1}\sigma_1^{\,k}z=\sigma_2^{-1}\sigma_1^{\,k}\sigma_2\in M$
for some $k\ne0$.  Consider the associated lifting loop $\ell$ with deck
displacement $g_\ell=\sigma_2^{-1}\sigma_1^{\,k}\sigma_2$.  Since
$\sigma_2\in Z$ maps $\D_f$ to itself and acts on $\Teich(\Sph,B)$ as a
deck transformation of $\pi$, the translate by $\sigma_2^{-1}$ of the
lifted path of $\ell$ starting at $\sigma_2\tau_0$ is a lifted path of
some loop $\ell'$ starting at $\tau_0$, remaining in $\D_f$, with deck
displacement
$$g_{\ell'}=\sigma_2^{-1}g_\ell\sigma_2
=\sigma_2^{-2}\sigma_1^{\,k}\sigma_2^{\,2}.$$
Now apply the winding argument of Step 4 of the proof of
Proposition~\ref{prop:input}.  All three elements $\sigma_1,\sigma_2,
\sigma_3$ are pure $v$--pushes (Lemma~\ref{lem:birman}), hence so is
$g_{\ell'}$; the winding of the $v$--strand of the braid $g_{\ell'}$ around
$c_1$ is computed on the abelianization of the point--pushing subgroup
$\Pi\cong\pi_1(\Sph\setminus A)\cong F_3$: writing $x_1,x_2,x_3$ for the
generators (pushes around $c_1,p_2,p_3$ respectively), we have
$\sigma_1=x_1^{2},\ \sigma_2=x_3^{2}$ (pushes around $p_3$) up to the
relation $x_1x_2x_3x_4=1$ of the four punctures, under which
$g_{\ell'}=x_3^{-2}x_1^{\,k}x_3^{2}$ has $x_1$--winding exactly $k\neq0$
and trivial $x_2$--winding.  A nonzero $x_1$--winding of the $v$--strand
can only be produced by loops with nonzero winding around the golden
cusps $a_\pm$ (Lemma~\ref{lem:monodromytable}: these are the only cusps at
which $v$ winds around $c_1$), contradicting
Proposition~\ref{prop:input}.  Hence \eqref{eq:separation} holds, and
Proposition~\ref{prop:cosets}(2) concludes.
\end{proof}

\begin{proof}[Proof of Theorem~\ref{thm:main}]
(a) is Theorem~\ref{thm:preserve}.  (b) is
Theorem~\ref{thm:infinite}.  The parenthetical statement about the curves
enclosing $\{p_2,v\}$ and $\{p_3,v\}$ follows from
Theorem~\ref{thm:preserve} and the same argument with the roles of the
indices permuted (using $\sigma_2$, resp.\ $\sigma_3$, as the conjugating
element $z$ and Proposition~\ref{prop:input} at the corresponding cusps).
\end{proof}

\section{Concluding remarks}

\begin{enumerate}
\item Theorem~\ref{thm:main} shows in particular that the deformation space
of a quadratic rational map with a critical point of period $4$ and the
second critical value off the cycle has \emph{infinitely many} connected
components, while its image in moduli is the (irreducible, rational)
stratum curve of \S\ref{s:stratum}.  Each component of $\D_f$ nevertheless
surjects (up to a finite set) onto the whole stratum: this follows from
Proposition~\ref{prop:image}(4) together with the coset description of
Proposition~\ref{prop:cosets}.
\item The mechanism has a transparent physical interpretation.  The stratum
carries four ``golden cusps'' $a_\pm,b_\pm$ at which the second critical
value collides with a point of the superattracting cycle; because
$v(t)=\bigl(\frac{t^2-t-1}{t^2+t-1}\bigr)^2$ is a square, the collision
has multiplicity two, so the vanishing cycle contributes a
\emph{square} Dehn twist.  The combinatorial class of the map changes by
this square twist when one encircles the cusp, so the square twist lies in
$Z$ but not in the monodromy group $M$ up to the relevant conjugacies; the
free subgroup $F\cong F_3$ of Proposition~\ref{prop:free} then guarantees
that no nonzero power of the twist is ever recaptured by $M$ after
conjugation.
\item The same framework applies verbatim to the curves enclosing
$\{p_2,v\}$ and $\{p_3,v\}$.  By contrast, as recorded in
\S\ref{s:preserve}, curves enclosing pairs of cycle points, or the pair
$\{p_1,v\}$, never satisfy the preservation criterion at any power: their
preimage structure forces a mismatch between the lifted multitwist and the
forgotten twist.  This ``asymmetry''---the special role of the critical
value $v=f(c_2)$ versus the cycle value $p_1=f(c_1)$---is a characteristic
feature of Epstein's deformation spaces for maps with a periodic critical
point.
\end{enumerate}

\begin{thebibliography}{99}

\bibitem{Birman} J.~S. Birman, \emph{Braids, links, and mapping class
groups}, Annals of Mathematics Studies, No. 82, Princeton University
Press, 1974.

\bibitem{BH} B. Branner and J.~H. Hubbard, \emph{The iteration of cubic
polynomials, Part I: The global topology of parameter space}, Acta Math.
\textbf{160} (1988), 143--206.

\bibitem{EpsteinTransversality} A. Epstein, \emph{Transversality principles
in holomorphic dynamics}, preprint; and \emph{Limits of quadratic rational
maps}, preprint, Institute for Mathematical Sciences, SUNY Stony Brook.

\bibitem{FM} B. Farb and D. Margalit, \emph{A primer on mapping class
groups}, Princeton Mathematical Series 49, Princeton University Press,
2012.

\bibitem{Hubbard} J.~H. Hubbard, \emph{Teichm\"uller theory and
applications to geometry, number theory, and dynamics, Vol.~1}, Matrix
Editions, Ithaca, 2006.

\bibitem{Kiwi} J. Kiwi, \emph{Wandering orbit portraits}, Trans.\ Amer.\
Math.\ Soc.\ \textbf{354} (2002), 1473--1485.

\bibitem{MSS} R. Ma\~n\'e, P. Sad, and D. Sullivan, \emph{On the dynamics
of rational maps}, Ann. Sci.\ \'Ec.\ Norm.\ Sup\'er.\ (4) \textbf{16}
(1983), 193--214.

\bibitem{McMullen} C.~T. McMullen, \emph{Complex dynamics and
renormalization}, Annals of Mathematics Studies 135, Princeton University
Press, 1994.

\bibitem{PilgrimTan} K.~M. Pilgrim and Tan Lei, \emph{Rational maps with
fixed critical points}, in: Complex dynamics and related topics (Y.~Jiang,
ed.), International Press, 2004; see also K. Pilgrim, \emph{Combinatorial
dynamics and dynamics of branched covers}, Cambridge University Press,
~2000.

\bibitem{Rees} M. Rees, \emph{A partial description of parameter space of
rational maps of degree two, Part I}, Acta Math.\ \textbf{168} (1992),
11--87.

\bibitem{Thurston} W.~P. Thurston, \emph{On the geometry and dynamics of
iterates of rational maps}, with an appendix by D. Schleicher, in: Complex
dynamics: families and friends (D. Schleicher, ed.), A K Peters, 2009,
3--137.

\end{thebibliography}

\end{document}
```

## Summary of the answer

**The curve.** Take $\gamma\subset(\S^2,B)$ to be any simple closed curve enclosing **exactly the two points $\{c_1, v\}$** — the critical point of period $4$ together with the other critical value — i.e. one side of $\gamma$ is a disk $D$ with $D\cap B=\{c_1,v\}$. (The curves enclosing $\{p_2,v\}$ or $\{p_3,v\}$ work equally well.)

**The power.** The **square** $T_\gamma^{\,2}$ preserves $\mathcal D_f$. The factor $2$ is forced twice over: (i) topologically, $f^{-1}(\gamma)$ is a *single* curve mapping onto $\gamma$ with degree $2$, and a Dehn twist lifts through a degree-$d$ cover only after raising to the $d$-th power; (ii) algebraically, in the normalized family $f_t(z)=\frac{tz^2-(t^2+t+1)z+t^2+t}{z^2-(t^2+t+1)z+t^2+t}$ of quadratic rational maps with cycle $0\to1\to\infty\to t\to0$, the second critical value is the *square* $v(t)=\bigl(\frac{t^2-t-1}{t^2+t-1}\bigr)^2$, so at the golden cusps $t=\frac{1\pm\sqrt5}{2}$ (where $v\to c_1$) one turn of the parameter produces two turns of $v$ around $c_1$ — precisely one unit of the twist $T_{\{c_1,v\}}^2$.

**Why the orbits on components are infinite.** The full preimage $f^{-1}(\gamma)$ is a curve enclosing $\{f^{\circ 3}(c_1), w_0, c_2\}$, which is peripheral in $(\mathbb S^2,A)$; hence the lifted twist is trivial on $\mathrm{Teich}(\mathbb S^2,A)$ and matches the trivial action of the forgotten twist — this is the preservation criterion $\Theta=\pi_1$ (Theorem 4.6 of the paper). Fibres of $\mathcal D_f$ over the base configuration identify with the subgroup $Z\le\mathrm{Mod}(\mathbb S^2,B)$ of marked self-conjugacy classes; $Z$ contains the free group $\langle T_{\{c_1,v\}}^2,T_{\{p_2,v\}}^2,T_{\{p_3,v\}}^2\rangle\cong F_3$ (squares of point pushes). Components of $\mathcal D_f$ are right cosets $Z/M$ of the monodromy group $M$, and $T_\gamma^2$ acts by left multiplication. Because the combinatorial class of the map necessarily changes when one loops around a golden cusp (the walls issuing from the cusps — proved via the invariant real “carousal” of the four postcritically finite limits $f_{\frac{1\pm\sqrt5}{2}}, f_{\frac{-1\pm\sqrt5}{2}}$ — separate the classes, and the sign of $c_2(t)-t$ flips across each cusp), no nontrivial power of $T_{\{c_1,v\}}^2$, nor any of its conjugates by $T_{\{p_2,v\}}^2$, lies in $M$. Hence the cosets $\sigma_1^k\sigma_2 M$, $k\in\mathbb Z$, are all distinct: the twist has infinite orbits on $\pi_0(\mathcal D_f)$, and $\mathcal D_f$ has infinitely many connected components.